\section{A simultaneous finite-place theta regularizer}
\label{sec:sr-finite-regularizer}

This section constructs the regularizer used to pass from matched
quotients to ordinary convergent theta integrals.  The construction uses
only the split Schr\"odinger model at one finite place.  Its finite support,
nonzero degree, and compatibility with a finite set of relations are
proved explicitly.

Let \(k\) be a nonarchimedean local field with valuation ring \(\mathcal O\),
uniformizer \(\pi\), and residue field \(\mathbb F_Q\).  Give \(k\) the
absolute value \(|\pi|=Q^{-1}\).  Let \(d,r\) be nonnegative integers with
\(d/2\ge r+1\), and let \(\nu:k^\times\to\mathbb C^\times\) be unitary.
On functions on \(M_{r,d}(k)\) use the split Weil action
\begin{equation}\label{eq:sr-regularizer-action}
 (T_g\varphi)(Y)=\nu(\det g)|\det g|^{d/2}\varphi(g^{\mathsf t}Y),
 \qquad g\in\operatorname{GL}_r(k).
\end{equation}
A row frame \(Y\in M_{r,d}(\mathcal O)\) is primitive if its reduction
modulo \(\pi\) has row rank \(r\).  Equivalently its rows are independent
over \(k\) and their \(\mathcal O\)-span is saturated in
\(\mathcal O^d\).  Indeed, independence of the reductions implies that
every \(c\in k^r\) with \(cY\in\mathcal O^d\) belongs to \(\mathcal O^r\):
otherwise multiply \(c\) by a power of \(\pi\) to obtain a vector with at
least one unit coordinate whose product with \(Y\) is zero modulo \(\pi\).
The converse follows by lifting a nonzero relation modulo \(\pi\) and
dividing it by \(\pi\).

\begin{srproposition}[An explicit primitive-frame operator]
\label{prop:sr-primitive-regularizer}
There is a finitely supported sum
\[
 Z=\sum_{g\in\operatorname{GL}_r(k)}z_g T_g
\]
such that
\[
 Z{\bf1}_{M_{r,d}(\mathcal O)}
   ={\bf1}_{\{\text{primitive row frames}\}}.
\]
Writing \(x=\nu(\pi)Q^{-d/2}\) and \(\deg Z=\sum_gz_g\), one has
\begin{equation}\label{eq:sr-regularizer-degree}
 \deg Z=\prod_{i=0}^{r-1}(1-Q^i x),\qquad
 |\deg Z|\ge\frac12,\qquad
 \sum_g|z_g|\le e^{1/2}.
\end{equation}
In particular \(Z_0=(\deg Z)^{-1}Z\) has degree one and coefficient
\(\ell^1\)-norm less than \(4\).  With additive Haar measure on
\(M_{r,d}(k)\), its operator norm on \(L^2\) is less than \(4\).
\end{srproposition}

\begin{proof}
For every subspace \(L\le\mathbb F_Q^r\) of dimension \(j\), choose a
lift of a basis of \(L\), extended to a matrix \(A_L\in
\operatorname{SL}_r(\mathcal O)\).  For \(L=\mathbb F_Q^r\) take
\(A_L=1\).  Such an extension exists for proper \(L\): extend the basis
to an invertible matrix and adjust one of the complementary rows to
make its determinant one, then lift an element of
\(\operatorname{SL}_r(\mathbb F_Q)\) to
\(\operatorname{SL}_r(\mathcal O)\).
Put
\[
 \gamma_L=\operatorname{diag}(\underbrace{\pi^{-1},\ldots,\pi^{-1}}_j,
                         1,\ldots,1)A_L,\qquad g_L=\gamma_L^{\mathsf t}.
\]
Then \(\det g_L=\pi^{-j}\), and
\[
 \gamma_LY\in M_{r,d}(\mathcal O)
 \quad\Longleftrightarrow\quad
 Y\in M_{r,d}(\mathcal O)\ \hbox{and}\ L\overline Y=0.
\]
For the reverse implication to integrality, use
\(\gamma_L^{-1}=A_L^{-1}\operatorname{diag}(\pi,\ldots,\pi,1,\ldots,1)\),
whose entries are integral.

The M\"obius function from zero to a \(j\)-dimensional subspace in the
subspace lattice is
\(\mu(L)=(-1)^jQ^{j(j-1)/2}\).  This follows inductively from
\(\sum_{L\le U}\mu(L)=0\) for every nonzero \(U\); equivalently it is the
coefficient identity in the finite \(Q\)-binomial theorem
\[
 \sum_{L\le\mathbb F_Q^r}
    (-1)^{\dim L}Q^{\dim L(\dim L-1)/2}X^{\dim L}
 =\prod_{i=0}^{r-1}(1-Q^iX).
\]
The identity itself follows by induction on \(r\), splitting the
subspaces according to their intersection with a fixed hyperplane; a
\(j\)-dimensional subspace not contained in that hyperplane is specified
by its \((j-1)\)-dimensional intersection and \(Q^{r-j}\) possible
lifts.

Take
\[
 Z=\sum_{L\le\mathbb F_Q^r}\mu(L)x^{\dim L}T_{g_L}.
\]
The scalar in \(T_{g_L}\) is \(x^{-j}\); hence, at an integral \(Y\),
the value of \(Z{\bf1}_{M_{r,d}(\mathcal O)}\) is
\(\sum_{L\le\ker(\overline Y)}\mu(L)\).  This is one precisely when
\(\overline Y\) has independent rows and zero otherwise.  At a
nonintegral \(Y\) every summand is zero, by the preceding equivalence.

The same \(Q\)-binomial identity gives the degree in
\eqref{eq:sr-regularizer-degree}.  Set \(t_i=Q^{i-d/2}\).  Since
\(d/2\ge r+1\),
\[
 \sum_{i=0}^{r-1}t_i
 \le \sum_{j=2}^{\infty}Q^{-j}
 =\frac{1}{Q(Q-1)}\le\frac12.
\]
Consequently
\[
 |\deg Z|\ge\prod_{i=0}^{r-1}(1-t_i)
 \ge1-\sum_{i=0}^{r-1}t_i\ge\frac12.
\]
Taking absolute values of the coefficients before combining repeated
\(g_L\)'s, the \(Q\)-binomial identity with \(-|x|\) gives
\[
 \sum_g|z_g|
 \le\sum_LQ^{j(j-1)/2}|x|^j
 =\prod_{i=0}^{r-1}(1+t_i)
 \le\exp\left(\sum_it_i\right)\le e^{1/2}.
\]
Each \(T_g\) is an \(L^2\) isometry: the change of variables
\(Y\mapsto g^{\mathsf t}Y\) has modulus \(|\det g|^d\), cancelling the
square of the scalar in \eqref{eq:sr-regularizer-action}.  The operator
bound follows from the triangle inequality.
\end{proof}

The construction is unchanged by adapted integral coordinates.  More
precisely, let \(\Lambda\) be an \(\mathcal O\)-lattice in \(k^d\), let
\(P\in\operatorname{GL}_r(k)\), and take the local vacuum
\[
 \varphi_0(Y)=c\,{\bf1}_{M_r(\Lambda)}(PY).
\]
Use the preceding proof in a basis of \(\Lambda\) and replace every
\(g\) by \(P^{\mathsf t}g(P^{\mathsf t})^{-1}\).  Conjugation preserves
determinants and coefficient sums.  The resulting operator sends
\(\varphi_0\) to \(c\) times the primitive-frame indicator at \(PY\),
and satisfies the same degree and coefficient bounds.  Primitive
frames form a compact open subset of the full-rank locus.  For example,
compactness follows because they are a closed subset of
\(M_{r,d}(\mathcal O)\), and openness follows from their description by
the nonvanishing of an \(r\)-minor modulo \(\pi\).

Now let \(E/F\) be a CM extension and consider a small theta pair
\[
 \dim_E V'=r=a',\qquad \dim_E W=d=a+a',\qquad a'<a,
\]
with \(d\) even.  Thus \(a\ge a'+2\) and \(d/2\ge r+1\), exactly the
range of Proposition~\ref{prop:sr-primitive-regularizer}.  At a split
finite place \(v\) of \(F\), the fixed integral Schr\"odinger model
identifies the action of \(U(V')(F_v)\) with
\eqref{eq:sr-regularizer-action}, possibly after the adapted coordinate
change just described.  This identification includes its scalar
\(\nu(\det g)|\det g|^{d/2}\); replacing it by a scalar-free action would
give the wrong degree.

\begin{srproposition}[One place for a finite collection]
\label{prop:sr-simultaneous-regularizer}
Let \(\Psi_1,\ldots,\Psi_N\) belong to the algebraic restricted tensor
product of the fixed small theta model, and let
\(u_1,\ldots,u_N\in U(V')(\mathbb A_{F,f})\).
One can choose a single split finite place \(v\), one operator \(Z\)
supported in \(U(V')(F_v)\), and finite adelic elements \(u_i^{\,v}\)
whose \(v\)-components are one, such that
\[
 \deg Z\ne0,\qquad
 Z\omega(u_i)\Psi_i=\omega(u_i^{\,v})Z\Psi_i\quad(1\le i\le N).
\]
Every \(Z\Psi_i\) has a finite sum of pure tensor representations whose
\(v\)-factors are supported in the full-rank locus.  Assume the fixed arithmetic pair is the nondegenerate unitary Weil
pair of Section~\ref{sec:theta-decay-proof}, with its rational
Schr\"odinger coordinates, standard additive character, and the
unitary reduction-theoretic cover, height comparisons, and finite
quotient volume stated there.  Then every such vector has an ordinary
absolutely convergent theta integral of moderate growth, by
\eqref{eq:theta-full-decay} and the coefficient proof in
Subsection~\ref{subsec:sr-theta-coefficient-proof}.
\end{srproposition}

\begin{proof}
Write each \(\Psi_i\) as a finite sum of pure tensors.  Outside a finite
set of places all these tensors have the same prescribed vacuum
factor.  Enlarge this finite set so that the local characters and
integral models are unramified and each \(u_i\) belongs to the
hyperspecial group fixing the vacuum at every remaining place.
There are infinitely many rational primes splitting completely in a
Galois closure of \(E/\mathbb Q\); choose one outside the finite set and
take a place \(v\) of \(F\) above it.  This is the only use here of the
standard splitting-prime existence theorem; no size estimate for \(v\)
is required.

Use Proposition~\ref{prop:sr-primitive-regularizer} in the common vacuum
coordinates at \(v\), and transport \(Z\) into \(U(V')(F_v)\).  Its
degree stays nonzero.  The tensor factorization shows simultaneously
that all \(Z\Psi_i\) have primitive-frame factors there.  Define
\(u_i^{\,v}\) by deleting the \(v\)-component of \(u_i\).
The deleted component acts trivially on every original vacuum factor,
so \(\omega(u_i)\Psi_i=\omega(u_i^{\,v})\Psi_i\).
The operator \(Z\) commutes with \(\omega(u_i^{\,v})\), since they are
supported at disjoint places.  This proves the asserted identity.
Full-rank local support is precisely the hypothesis of
\eqref{eq:theta-full-decay}.  Under the fixed-pair and reduction
hypotheses in the statement, Subsection~\ref{subsec:sr-theta-coefficient-proof}
proves the coefficient majorant used by that estimate; the explicit
packing and minor argument in Section~\ref{sec:theta-decay-proof}
then proves convergence and moderate growth.
\end{proof}

The preceding proposition makes the quotient argument precise.
Let \(G=U(V')\), fix any Haar measure \(dg\) and a measurable fundamental
domain \(\mathcal F\) for \(G(F)\backslash G(\mathbb A_F)\), and write
\[
 I(\Psi)(h)=\int_{\mathcal F}\Theta_\Psi(g,h)\,dg
\]
whenever this integral is absolutely convergent.  The choice of Haar
measure is fixed throughout.  The theta transformation law and
rational left invariance imply
\begin{equation}\label{eq:sr-ordinary-translate}
 I(\omega(u)\Psi)=I(\Psi)
\end{equation}
whenever \(I(\Psi)\) exists.  Indeed
\(\Theta_{\omega(u)\Psi}(g,h)=\Theta_\Psi(gu,h)\);
right translation preserves quotient measure because a unitary group
is unimodular.  The two fundamental domains give the same integral
by rational left invariance.  In particular, finite sums satisfy
\begin{equation}\label{eq:sr-ordinary-degree}
 I(Z\Psi)=(\deg Z)I(\Psi)
\end{equation}
for every integrable \(\Psi\), regardless of the support of \(Z\).

\begin{srcorollary}[Integrable finite relations vanish]
\label{cor:sr-integrable-relations}
Suppose
\[
 \Psi=\sum_{i=1}^{N}b_i\bigl(\omega(u_i)\Psi_i-\Psi_i\bigr)
\]
belongs to the algebraic model, and its theta integral is absolutely
convergent.  Then \(I(\Psi)=0\).
\end{srcorollary}

\begin{proof}
Apply Proposition~\ref{prop:sr-simultaneous-regularizer} to the entire
finite collection.  Each \(Z\Psi_i\) is integrable and
\[
 Z\Psi=\sum_i b_i
 \bigl(\omega(u_i^{\,v})Z\Psi_i-Z\Psi_i\bigr).
\]
Every difference on the right has zero integral by
\eqref{eq:sr-ordinary-translate}; finite additivity is justified by
absolute convergence.  Thus \(I(Z\Psi)=0\).
On the other hand the assumed integrability of \(\Psi\) gives
\eqref{eq:sr-ordinary-degree}.  Divide by the nonzero degree.
\end{proof}

Consequently, if two small vectors have the same image in the finite
coinvariant quotient, both have convergent theta integrals, and their
difference is a finite relation of the displayed form, their ordinary
integrals agree.  If a representative \(\Psi_0\) in that class has
already been matched to a big vector \(\Phi\) with
\(I_V(\Phi)=cI_{V'}(\Psi_0)\), the same equality holds for every
integrable representative \(\Psi\) in the class.  There is no additional
volume scalar in this passage: the argument uses one Haar measure
throughout, and scaling it scales every small integral by the same
amount.

This constructs the finite-place regularizer and proves the passage
from a coinvariant class to any convergent representative.  The scalar
\(c\), its nonvanishing, the placewise matching maps, and the universal
character-independent rank-recursion signs still require their
separate theta-comparison proofs; they are not consequences of the
regularizer.  In the frozen source the construction here occurs in
\srcid{ExistsTranslatesEqPrimFrameLat},
\srcid{ExistsGoodPlaceTranslateCombinationLatticeIndica},
\srcid{ExistsCombinationMemRegularSubspace},
\srcid{TwistRegulariserIsGood}, and
\srcid{ThetaIntegralsAgreeOfIntegrableKernel}.

\subsection{The coefficient sum and its Schwartz majorant}
\label{subsec:sr-theta-coefficient-proof}

We give the coefficient calculation used in
\eqref{eq:theta-lattice-majorant}.
Let \(A=\mathbb A_F\), choose the standard nontrivial additive
character \(\psi:A/F\to\mathbb C^\times\), and give the additive
spaces self-dual Haar measure.  Use a rational polarization of
the symplectic space of the specified Weil pair, so its
Schr\"odinger space is \(\mathcal S(A^N)\) and its rational
Heisenberg coordinates are \(F^N\oplus F^N\).  In these coordinates
\[
 [S(a,b)f](u)=
 \psi\bigl(b\cdot u+\tfrac12 a\cdot b\bigr)f(u+a),
 \qquad Pf(u)=f(-u).
\]
The rational polarization identifies these coordinates with the
rational matrix space used in Section~\ref{sec:theta-decay-proof};
the action of \((g,h)\) on that space is
\(x\mapsto gx h^{\mathsf t}\), with the fixed changes of bases
included in the model.

For \(\Phi\in\mathcal S(A^N)\) set
\[
 c_\Phi(a,b)=
 \int_{A^N}\psi\bigl(b\cdot u+\tfrac12a\cdot b\bigr)
          \Phi(u+a)\overline{\Phi(-u)}\,du
 =\langle S(a,b)\Phi,P\Phi\rangle .
\]
Then
\begin{equation}\label{eq:sr-theta-square-coefficients}
 \left|\sum_{x\in F^N}\Phi(x)\right|^2
       =\sum_{a,b\in F^N}c_\Phi(a,b).
\end{equation}
All series in this identity converge absolutely.

Here are the analytic details justifying that assertion and its
uniform use with the infinite Schwartz variable.
At the archimedean places, \(c_f\) is a Schwartz function of \((a,b)\).
For every desired weight exponent its Schwartz seminorm is at most
a constant times \(p(f)^2\), for a continuous Schwartz seminorm \(p\).
To verify this, differentiate its integral in \(a,b\).
A derivative in \(b\) inserts a polynomial in \(u+a/2\);
a derivative in \(a\) differentiates \(f(u+a)\) and the phase.
To bound a monomial in \(b\), move derivatives of the phase onto
the two Schwartz factors by integration by parts.
Write each occurrence of \(a\) as \((u+a)-u\) and expand.
Every resulting integrand is a product of weighted derivatives of
\(f(u+a)\) and \(f(-u)\).
Bound the weighted derivatives by their Schwartz seminorms times
\((1+|u+a|)^{-m}(1+|u|)^{-m}\), with \(m\) larger than the real
dimension.  Its integral is bounded independently of \(a\).
This bounds every polynomially weighted derivative of \(c_f\).
Only finitely many derivatives and weights are needed for a fixed
bound, so their maximum, enlarged by a fixed constant, is the
continuous seminorm \(p\).
The same proof applies to the product of the real and complex
archimedean coordinates after the fixed linear change making the
character the standard Fourier character.

At a finite place, let \(\phi\) be locally constant with compact
support \(K\), and choose a compact open additive subgroup \(U\)
under which \(\phi\) is invariant.
Nonvanishing of the integrand implies \(a\in K+K\).
Translating the integration variable by \(t\in U\) multiplies
the coefficient by \(\psi(b\cdot t)\); hence it is zero unless
\(b\in U^\perp\).  The annihilator \(U^\perp\) is compact open.
The coefficient is locally constant, and thus is again
Schwartz--Bruhat, with compact support in
\((K+K)\times U^\perp\).
For a fixed finite adelic factor \(\phi_f\) this proves a bound
\[
 |c_{f\otimes\phi_f}(w)|
 \le p_{\phi_f,L}(f)^2\,
       {\bf1}_{C_f}(w_f)(1+|w|_\infty)^{-L}
 \qquad(w\in A^{2N})
\]
for every \(L\), some fixed compact \(C_f\), and a continuous
seminorm \(p_{\phi_f,L}\).
Restricted tensor products are handled by multiplication of the
local formulas; their vacuum factors agree outside finitely many
places.  Finite linear combinations follow by polarization and the
triangle inequality.

It remains to check \eqref{eq:sr-theta-square-coefficients}.
For fixed rational \(a\), apply additive adelic Poisson summation
in \(u\) to the Schwartz function
\(\Phi(u+a)\overline{\Phi(-u)}\).
Because \(a,b\) are rational, \(\psi(a\cdot b/2)=1\).
Self-dual measure for the standard character has
\(\operatorname{vol}(A^N/F^N)=1\), so there is no Poisson volume
factor, and the sum over rational \(b\) of its Fourier transform is
\[
 \sum_{u\in F^N}\Phi(u+a)\overline{\Phi(-u)}.
\]
Summing over rational \(a\), the bijection
\((a,u)\mapsto(u+a,-u)\) of \(F^N\times F^N\) gives the product
of the two theta sums in
\eqref{eq:sr-theta-square-coefficients}.
The coefficient sum is absolutely convergent by its proved Schwartz
bound.  The double sum of the products is absolutely convergent
because each theta sum of the absolute value of a Schwartz function
converges.  These facts justify both Poisson summation and the
rearrangement.

The unitary Weil action of the pair preserves the Schwartz space,
has
\(\omega(g,h)S(w)\omega(g,h)^{-1}=S((g,h)w)\), and commutes with
\(P\).  These assertions follow directly on the Schr\"odinger
generators: a linear change has the square-root modulus making it
unitary, multiplication by a quadratic character has modulus one
and is even, and the self-dual Fourier transform is unitary and
commutes with reflection.  Thus
\[
 c_{\omega(g,h)\Phi}(w)=c_\Phi((g,h)^{-1}w).
\]
Apply \eqref{eq:sr-theta-square-coefficients} to
\(\omega(g,h)\Phi\), take absolute values, and apply the preceding
Schwartz bound.  In the matrix coordinates of the fixed pair,
this is precisely
\[
 |\Theta_{f\otimes\phi_f}(g,h)|^2
 \le p_{\phi_f,L}(f)^2
     \sum_{x}
       {\bf1}_{C_f}\bigl((g^{-1}x(h^{-1})^{\mathsf t})_f\bigr)
       (1+|g^{-1}x(h^{-1})^{\mathsf t}|_\infty)^{-L}.
\]
For full-rank local data, apply the explicit compact-coset
Schr\"odinger separation in
Section~\ref{sec:theta-decay-proof} to the local factor before
this calculation.  On each piece the coefficient vanishes at
every rank-deficient local matrix.  Rational rank-deficient
matrices remain rank deficient under the two invertible group
matrices.  The sum can therefore be restricted to full row rank.
For a finite sum of pieces, take square roots of the separate
bounds, use the union of their finite compact supports, and add
their seminorms with the absolute values of their coefficients.
This proves \eqref{eq:theta-lattice-majorant} in both of its forms.

Consequently the lattice packing, full-rank minor estimate, and
two-chamber calculation in
Section~\ref{sec:theta-decay-proof} prove
\eqref{eq:theta-full-decay} for this fixed arithmetic pair.
The remaining standard reduction-theoretic assumptions are exactly
the unitary Siegel-set cover, its height comparisons, and finite
Haar volume of \(U_X(F)\backslash U_X(\mathbb A_F)\), as stated
there.  They apply to the nondegenerate Hermitian space of the
fixed pair; its unitary group is reductive and its center is the
norm-one torus of the CM extension, which is anisotropic over
\(F\).  At \(N_0=0\), \eqref{eq:theta-full-decay} is uniform in
\(g\); integration over this finite-volume quotient gives ordinary
absolute convergence.  Its polynomial bound in \(h\) gives moderate
growth of the resulting integral.  None of these assertions
constructs the arithmetic pair, nor asserts uniformity of its
Schwartz seminorm constants across varying fields.
